\documentclass[11pt,a4paper,oneside]{book}
\usepackage[utf8]{inputenc}
\usepackage[turkish]{babel}
\usepackage{amsmath,amssymb,amsfonts,amsthm}
\usepackage{geometry}
\geometry{a4paper, margin=2.2cm, top=2.5cm, bottom=2.5cm}
\usepackage{tcolorbox}
\tcbuselibrary{breakable,skins}
\usepackage{tikz}
\usepackage{pgfplots}
\pgfplotsset{compat=1.18}
\usetikzlibrary{babel,arrows.meta,positioning,calc,shapes.geometric}
\usepackage{hyperref}
\usepackage{booktabs}
\usepackage{tabularx}
\usepackage{centernot}
\usepackage{caption}
\usepackage{fancyhdr}
\usepackage{etoolbox}
\usepackage{docmute}

\hypersetup{
    colorlinks=true,
    linkcolor=primaryblue,
    citecolor=defgreen,
    urlcolor=primaryblue,
    bookmarksnumbered=true,
    pdfstartview={XYZ null null 1.00},
    pdftitle={Istatistik Teorisi Ders Kitabi},
    pdfauthor={Istatistik Teorisi Notlari}
}

%% Renk Paleti
\definecolor{primaryblue}{RGB}{25, 60, 120}
\definecolor{secondaryblue}{RGB}{235, 242, 250}
\definecolor{defgreen}{RGB}{34, 112, 60}
\definecolor{lightgreen}{RGB}{236, 247, 240}
\definecolor{warnred}{RGB}{168, 38, 38}
\definecolor{lightred}{RGB}{253, 242, 242}
\definecolor{examyellow}{RGB}{180, 115, 20}
\definecolor{lightyellow}{RGB}{254, 250, 238}

% Aliases for Hafta 01-06 color names
\definecolor{mainblue}{RGB}{25, 60, 120}
\definecolor{okgreen}{RGB}{34, 112, 60}
\definecolor{alertred}{RGB}{168, 38, 38}
\definecolor{purp}{RGB}{128, 0, 128}
\definecolor{boxbg}{RGB}{245, 248, 252}
\definecolor{tealblue}{RGB}{0, 120, 130}

%% Kutu Tanımları
\tcbset{
    commonbox/.style={
        breakable,
        enhanced,
        arc=2mm,
        boxrule=1pt,
        left=4mm, right=4mm, top=3mm, bottom=3mm
    }
}

\newtcolorbox[auto counter, number within=chapter]{tanim}[2][]{
    commonbox,
    colback=lightgreen,
    colframe=defgreen,
    coltitle=white,
    fonttitle=\bfseries,
    title={Tan\i m \thetcbcounter\notblank{#2}{ --- #2}{}},
    #1
}

\newtcolorbox[auto counter, number within=chapter]{teorem}[2][]{
    commonbox,
    colback=secondaryblue,
    colframe=primaryblue,
    coltitle=white,
    fonttitle=\bfseries,
    title={Teorem \thetcbcounter\notblank{#2}{ --- #2}{}},
    #1
}

\newtcolorbox[auto counter, number within=chapter]{ornek}[2][]{
    commonbox,
    colback=lightyellow,
    colframe=examyellow,
    coltitle=white,
    fonttitle=\bfseries,
    title={\"{O}rnek \thetcbcounter\notblank{#2}{ --- #2}{}},
    #1
}

\newtcolorbox[auto counter, number within=chapter]{uyari}[1][]{
    commonbox,
    colback=lightred,
    colframe=warnred,
    coltitle=white,
    fonttitle=\bfseries,
    title={Hata Tespiti / Uyar\i\ \thetcbcounter\notblank{#1}{ --- #1}{}},
}

\newtcolorbox[auto counter, number within=chapter]{odev}[2][]{
    commonbox,
    colback=lightgreen!60!secondaryblue,
    colframe=defgreen!80!black,
    coltitle=white,
    fonttitle=\bfseries,
    title={\"{O}dev \c{C}\"{o}z\"{u}m\"{u} \thetcbcounter\notblank{#2}{ --- #2}{}},
    #1
}

\newtcolorbox[auto counter, number within=chapter]{sezgi}[2][]{
    commonbox,
    colback=tealblue!5!white,
    colframe=tealblue,
    coltitle=white,
    fonttitle=\bfseries,
    title={Sezgisel Bak\i\c{s} ve Basit Anlat\i m\notblank{#2}{ --- #2}{}},
    #1
}

%% Sayfa Başlıkları
\pagestyle{fancy}
\fancyhf{}
\fancyhead[L]{\small\bfseries\color{primaryblue}\leftmark}
\fancyhead[R]{\small\bfseries\thepage}
\renewcommand{\headrulewidth}{0.6pt}
\renewcommand{\chaptermark}[1]{\markboth{\thechapter.\ #1}{}}

\begin{document}
\shorthandoff{=}

%% Kapak
\begin{titlepage}
\centering
\vspace*{1.5cm}
{\Huge\bfseries\color{primaryblue} \.{I}STAT\.{I}ST\.{I}K TEOR\.{I}S\.{I} \\[0.4em] \LARGE DERS K\.{I}TABI}\\[1.5cm]
{\Large\bfseries Kapsaml\i\ Teori, Matematiksel \.{I}spatlar, Hata D\"{u}zeltmeleri,\\[0.3em]
\"{O}zg\"{u}n TikZ Grafikleri ve Ayd\i nlat\i lm\i \c{s} \"{O}dev \c{C}\"{o}z\"{u}mleri}\\[2cm]

\begin{tcolorbox}[enhanced, colback=secondaryblue, colframe=primaryblue, width=0.88\textwidth, arc=2mm, boxrule=1pt]
\centering\large
\textbf{Kaynak Materyal:}\\
45 Sayfal\i k Orijinal El Yaz\i s\i\ Ders Notlar\i\ (2023 Bahar D\"{o}nemi)\\
\vspace{0.4em}
\normalsize
01--05.pdf $\cdot$ 06--11.pdf $\cdot$ 12--18.pdf $\cdot$ 19--24.pdf $\cdot$ 25--36.pdf $\cdot$ 37--42.pdf $\cdot$ 43--45.pdf
\end{tcolorbox}

\vfill
{\large \textbf{Haftal\i k Ders Notlar\i ndan Dinamik Olarak Ba\u{g}lanan Ana Metin}}\\[0.5em]
{\large 2026}
\vspace*{1cm}
\end{titlepage}

\frontmatter

\chapter*{\"{O}ns\"{o}z}
Bu ders kitab\i, Lisans ve Lisans\"{u}st\"{u} d\"{u}zey \textbf{\.{I}statistik Teorisi} dersi kapsam\i ndaki t\"{u}m temel konular\i\ tek bir kaynakta toplamak amac\i yla haz\i rlanm\i\c{s}t\i r. 

Klasik teorik anlat\i mlar\i n aksine, bu kitapta her matematiksel ispat\i n ard\i ndaki \textbf{sezgisel d\"{u}\c{s}\"{u}nce}, g\"{u}nl\"{u}k hayat analojileri, orijinal el yaz\i s\i\ notlardaki \textbf{hata d\"{u}zeltmeleri} ve \textbf{vekt\"{o}rel TikZ grafikleri} titizlikle i\c{s}lenmi\c{s}tir.

Kitap, \texttt{haftalik\_notlar/} alt\i ndaki m\"{o}d\"{u}ler haftal\i k not dosyalar\i n\i\ dinamik olarak (\texttt{\textbackslash docmute} ve \texttt{\textbackslash input}) \c{c}ekerek olu\c{s}turulmaktad\i r. B\"{o}ylece haftal\i k notlarda yap\i lan herhangi bir iyile\c{s}tirme veya d\"{u}zeltme, do\u{g}rudan bu kitaba yans\i maktad\i r.

\tableofcontents

\mainmatter

%% Haftalık alt belgelerin başlık/özet davranışları
\renewcommand{\maketitle}{}
\newenvironment{abstract}{%
  \begin{tcolorbox}[
    enhanced,
    colback=secondaryblue!50,
    colframe=primaryblue!70,
    arc=2mm,
    boxrule=0.8pt,
    left=4mm, right=4mm, top=3mm, bottom=3mm
  ]
  \small\textbf{Haftal\i k Konu ve B\"{o}l\"{u}m \"{O}zeti:}\par\smallskip\ignorespaces
}{%
  \end{tcolorbox}\vspace{1em}%
}

%% BÖLÜM 1
\chapter{Momentler, MGF ve Temel E\c{s}itsizlikler}
\input{../haftalik_notlar/Hafta_01.tex}

%% BÖLÜM 2
\chapter{Rastgele De\u{g}i\c{s}kenlerin D\"{o}n\"{u}\c{s}\"{u}mleri ve \"{O}zel Da\u{g}\i l\i mlar}
\input{../haftalik_notlar/Hafta_02.tex}

%% BÖLÜM 3
\chapter{Çok Boyutlu D\"{o}n\"{u}\c{s}\"{u}mler ve Da\u{g}\i l\i mlar ($t$, $F$, Beta)}
\input{../haftalik_notlar/Hafta_03.tex}

%% BÖLÜM 4
\chapter{S\i ra \.{I}statistikleri ve Asimptotik Limit Da\u{g}\i l\i mlar\i}
\input{../haftalik_notlar/Hafta_04.tex}

%% BÖLÜM 5
\chapter{Stokastik Yak\i nsama ve Asimptotik Teori}
\input{../haftalik_notlar/Hafta_05.tex}

%% BÖLÜM 6
\chapter{Nokta Tahmini, Sapmas\i zl\i k ve Tutarl\i l\i k}
\input{../haftalik_notlar/Hafta_07.tex}

%% BÖLÜM 7
\chapter{En \c{C}ok Olabilirlik Tahmini (MLE) ve CRLB Alt S\i n\i r\i}
\input{../haftalik_notlar/Hafta_06.tex}

\end{document}
